\section{A matrix energy and its Hessian}\label{sec:matrix}

We begin with the finite-dimensional inequalities that control the
derivatives of our energy.  Their constants are independent of the
matrix dimensions.  Throughout this section, matrices act on
finite-dimensional complex Hilbert spaces, and inequalities between
Hermitian matrices mean inequalities of quadratic forms.

For a positive semidefinite matrix $P$, set
\[
 \Phi(P)=\tr(P^{3/2}),
 \qquad \supp P=(\ker P)^\perp.
\]
A matrix is \emph{supported on $P$} if it vanishes on $\ker P$ on
both sides.  We use the Hilbert--Schmidt inner product
$\langle K,L\rangle_{\HS}=\tr(K^*L)$.
For a Hermitian matrix $A$, write
$\mathsf L_A(X)=AX$ and $\mathsf R_A(X)=XA$.
On matrices acting on $\supp P$, the operator
$\mathsf L_{\sqrt P}+\mathsf R_{\sqrt P}$ is positive definite.
For Hermitian $K,L$ supported on $P$, define
\[
 \mathcal B_P(K,L)
 =\frac32\operatorname{Re}
   \big\langle K,
    (\mathsf L_{\sqrt P}+\mathsf R_{\sqrt P})^{-1}L
   \big\rangle_{\HS},
 \qquad b_P(K)=\mathcal B_P(K,K).
\]
If $P=0$, the only supported directions are zero, and these expressions
are defined to be zero.

The square-root derivative used here is a case of the classical
divided-difference calculus for Hermitian matrix
functions~\cite[Section~1, Equation~(1)]{Daletskii1957}.
We give a direct inverse-function proof to keep the support conventions explicit.

\begin{lemma}[Energy derivatives]\label{lem:hessian}
On the cone of matrices positive definite on a fixed support,
$\Phi$ is smooth.  For Hermitian directions $K,L$ on that support,
\[
 d\Phi_P(K)=\frac32\tr(\sqrt P K),
 \qquad
 d^2\Phi_P(K,L)=\mathcal B_P(K,L).
\]
If $p_i>0$ are the eigenvalues on the support, then in an orthonormal
eigenbasis
\begin{equation}\label{eq:hessian-metric}
 b_P(K)=\frac32\sum_{i,j}
          \frac{|K_{ij}|^2}{\sqrt{p_i}+\sqrt{p_j}}.
\end{equation}
In particular, $\mathcal B_P$ is a positive definite symmetric bilinear
form on the real space of supported Hermitian directions, and
\[
 |\mathcal B_P(K,L)|\le b_P(K)^{1/2}b_P(L)^{1/2}.
\]
\end{lemma}

\begin{proof}
Restrict to the support and put $A=\sqrt P$.
The derivative of the squaring map at $A$ is $X\mapsto AX+XA$.
In an eigenbasis of $A$, this map multiplies its $(i,j)$ entry by
$\sqrt{p_i}+\sqrt{p_j}>0$.  The inverse function theorem therefore
gives a smooth positive square root.  Its derivative $X_K$ in
direction $K$ satisfies
\[
 AX_K+X_KA=K.
\]
Taking a trace after multiplication by $A$ yields
$\tr(PX_K)=\tfrac12\tr(AK)$.  Differentiation of
$\Phi(P)=\tr(P\sqrt P)$ now gives
\[
 d\Phi_P(K)=\tr(KA)+\tr(PX_K)=\frac32\tr(AK).
\]
Differentiating this identity in direction $L$ proves the Hessian
formula.  The entrywise formula follows by diagonalizing $A$.
It also proves symmetry, positivity, and Cauchy--Schwarz on the
real space of Hermitian directions.
\end{proof}

The energy and its Hessian are unchanged by simultaneous unitary
conjugation of their arguments.  They are also unchanged if all
arguments are extended by zero to a larger Hilbert space: the
positive eigenvalues and the supported entries in
\eqref{eq:hessian-metric} stay the same.  We will use these
observations to compare Grams acting on different subspaces.

\Needspace{3\baselineskip}
\begin{lemma}[Order comparison]\label{lem:order}
Let $P,Q$ be positive semidefinite matrices on the same Hilbert space,
let $a\ge1$, and suppose $P\le aQ$.  Every Hermitian $K$ supported
on $P$ is supported on $Q$, and
\begin{equation}\label{eq:matrix-order}
 b_P(K)\ge a^{-1/2}b_Q(K).
\end{equation}
\end{lemma}

\begin{proof}
If $x\in\ker Q$, then $0\le\langle Px,x\rangle
\le a\langle Qx,x\rangle=0$, so $Px=0$.
Thus $\ker Q\subseteq\ker P$, proving the support assertion.

We recall the two operator-order facts needed for the comparison.
Inversion reverses order on positive definite matrices: for
$0<X\le Y$, conjugate by $X^{-1/2}$, invert the resulting matrix
at least $I$, and conjugate back.  Moreover,
\[
 \sqrt X=\frac1\pi\int_0^\infty
       X(X+tI)^{-1}\,\frac{\dd t}{\sqrt t}.
\]
This norm-convergent formula follows from the scalar integral by
spectral calculus.  Since
$X(X+tI)^{-1}=I-t(X+tI)^{-1}$, inversion order proves that the
positive square root preserves order. This is the classical
Loewner--Heinz square-root monotonicity; see, for example,
Kwong~\cite[Theorem~1 and Section~2]{Kwong1975}. The supported
Hessian comparison below is derived from these order facts here.

For $\eta>0$, the assumptions give
\[
 P+\eta I\le a(Q+\eta I),
 \qquad
 \sqrt{P+\eta I}\le\sqrt a\,\sqrt{Q+\eta I}.
\]
If $A\le B$, then
$\mathsf L_A+\mathsf R_A\le\mathsf L_B+\mathsf R_B$ as
selfadjoint operators on Hilbert--Schmidt space.  Indeed, left and
right multiplication by the positive matrix $B-A$ have the
nonnegative quadratic forms
$\tr(X^*(B-A)X)$ and $\tr(X^*X(B-A))$.
Inverting the resulting positive definite operators gives
\[
 b_{P+\eta I}(K)\ge a^{-1/2}b_{Q+\eta I}(K).
\]
Finally let $\eta\downarrow0$.
Formula~\eqref{eq:hessian-metric}, applied in a full eigenbasis,
shows that both limits are the asserted supported costs: every
entry of $K$ meeting the kernel of the respective matrix is zero.
\end{proof}

We will apply the Hessian to Gram matrices with an insertion between
the two factors.  The following bound controls such derivatives by
the energy itself.

\begin{lemma}[Inserted Gram bound]\label{lem:insertion}
Let $R:E\to F$ be a linear map between finite-dimensional complex
Hilbert spaces, and let $L:F\to F$ be Hermitian.  The matrix $R^*LR$
is supported on $R^*R$, and
\begin{equation}\label{eq:insertion-bound}
 b_{R^*R}(R^*LR)
 \le\frac34\lVert L\rVert_{\op}^2\Phi(R^*R).
\end{equation}
\end{lemma}

\begin{proof}
The insertion annihilates $\ker R=\ker(R^*R)$ and is Hermitian,
so it is supported on the Gram.  The assertion is immediate if
$R=0$.  Otherwise take a singular-value decomposition, with positive
singular values $\rho_i$, and let $C$ be the compression of $L$ to
the corresponding left singular vectors.  Then
$\lVert C\rVert_{\op}\le\lVert L\rVert_{\op}$ and
\[
 b_{R^*R}(R^*LR)
 =\frac32\sum_{i,j}
   \frac{\rho_i^2\rho_j^2}{\rho_i+\rho_j}|C_{ij}|^2.
\]
For $x,y>0$, the arithmetic--geometric mean inequality gives
\[
 \frac{x^2y^2}{x+y}
 \le\frac{(xy)^{3/2}}2
 \le\frac{x^3+y^3}4.
\]
Each row and each column of $C$ has squared Euclidean norm at most
$\lVert L\rVert_{\op}^2$.  Consequently
\[
 b_{R^*R}(R^*LR)
 \le\frac38\sum_{i,j}(\rho_i^3+\rho_j^3)|C_{ij}|^2
 \le\frac34\lVert L\rVert_{\op}^2\sum_i\rho_i^3,
\]
which is \eqref{eq:insertion-bound}.
\end{proof}
